



\documentclass[11pt]{article}

\usepackage[a4paper,margin=25mm]{geometry}
\usepackage[T1]{fontenc}
\usepackage{amsmath}
\usepackage{booktabs}
\usepackage{graphicx}
\usepackage{longtable}
\usepackage{listings}
\usepackage{xcolor}
\usepackage{hyperref}
\usepackage{pdfpages}
\usepackage{titling}
\usepackage{fancyhdr}

\hypersetup{
  colorlinks=true,
  linkcolor=blue,
  urlcolor=blue,
  citecolor=blue
}
\lstset{
  basicstyle=\ttfamily\small,
  breaklines=true,
  columns=fullflexible,
  keepspaces=true,
  showstringspaces=false
}

% Move title upwards
\setlength{\droptitle}{-4.5cm}

\title{\color{white}
\hspace*{3cm}\LARGE\bfseries LaTeX for Physics\\
\hspace*{4cm}Reports, Equations, Figures, and References}
\author{\color{white}\hspace*{12cm}UCC Physics Society}
\date{\color{white}\hspace*{12.5cm}\today}

% Creates a header/footer
\pagestyle{fancy}
\fancyhf{}

% Header
\lhead{\includegraphics[height=40pt]{physoc.jpg}}
\rhead{Git and GitHub for Physics}

% Header and footer lines
\renewcommand{\headrulewidth}{0.4pt}
\renewcommand{\footrulewidth}{0.4pt}

% Footer
\cfoot{Page \thepage}

% Prevent header height warnings
\setlength{\headheight}{45pt}

\begin{document}

% PDF as background of the first page
\AddToShipoutPictureBG*{%
  \AtPageLowerLeft{%
    \includegraphics[
      width=\paperwidth,
      height=\paperheight
    ]{Front.pdf}%
  }%
}

\maketitle

{\hypersetup{linkcolor=white}
\color{white}
\tableofcontents
}

\newpage

\section{One shared document-writing guide}
LaTeX turns plain-text source into a typeset document. Edit \texttt{.tex},
not the generated PDF. The syntax is the same on Windows, macOS, and Linux;
installation is in the OS guides. Use LaTeX for equations, consistent structure,
cross-references, and bibliographies rather than arranging every page manually.
For an official introduction and reference material, see the
\href{https://www.overleaf.com/learn/latex/Learn_LaTeX_in_30_minutes}{Overleaf
LaTeX introduction} and the
\href{https://www.latex-project.org/help/documentation/}{LaTeX Project
documentation}.

\section{Overleaf or local TeX?}
\begin{description}
\item[Overleaf:] Browser-based compilation and collaboration without installing
a local TeX distribution. Requires network/account; check sharing settings.
\item[Local editor and TeX:] Offline builds and Git-friendly source, with
separate editor and TeX installations.
\end{description}
For a new Overleaf project, use New Project > Blank Project and follow its
project instructions. Do not upload confidential research without permission.
For local work, build this guide from the repository root:
\begin{lstlisting}
latexmk -cd -pdf -interaction=nonstopmode -halt-on-error LaTeX/LaTeX.tex
\end{lstlisting}
The \texttt{latexmk} command runs the necessary LaTeX and BibTeX passes.

\section{Guide and example files}
This guide is self-contained: copy its entire source into a new \texttt{.tex}
file or an Overleaf blank project and compile it. No other workshop files
are required. A local build needs a TeX distribution with the packages listed
at the top of the source.
\begin{description}
\item[\texttt{LaTeX.tex / LaTeX.pdf}:] This guide and its compiled PDF.
\item[\texttt{ref.bib}:] Optional example bibliography for practising BibTeX.
\item[\texttt{LaTeX\_meme.png}:] Optional image for practising image inclusion.
\end{description}
Files such as \texttt{.aux}, \texttt{.log}, \texttt{.toc}, \texttt{.out},
\texttt{.fls}, and \texttt{.fdb\_latexmk} are generated build state.
\texttt{.bbl} is generated bibliography content; \texttt{.synctex.gz} links
source and PDF positions. Do not edit these by hand.

\section{Text, structure, and lists}
\begin{lstlisting}
\section{Method}
Normal text, \textbf{bold}, and \emph{emphasis}.

A blank line starts a new paragraph.
\begin{itemize}
  \item First point.
  \item Second point.
\end{itemize}
\end{lstlisting}
Commands start with a backslash; required arguments use braces.
Environments have matching begin/end commands.
A percent sign starts a source comment. Literal special characters include
\verb|\%|, \verb|\_|, \verb|\&|, and \verb|\#|.
Do not make paragraphs with repeated \verb|\\| or align equations using spaces.
For a local size change, group it: \verb|{\large Larger text}|.

\subsection{Formatting, alignment, and font sizes}
The original tutorial's formatting examples can be written as:
\begin{lstlisting}
\textbf{Bold text}, \textit{italic text}, \underline{underlined text}.
\begin{center}
  Centred text.
\end{center}
\begin{flushright}
  Right-aligned text.
\end{flushright}
\begin{flushleft}
  Left-aligned text.
\end{flushleft}
{\Huge Very large text}\par
{\LARGE Larger text}\par
{\Large Large text}\par
{\normalsize Normal text}\par
{\small Small text}\par
{\footnotesize Footnote-sized text}\par
{\scriptsize Script-sized text}\par
{\tiny Tiny text}\par
\end{lstlisting}
Braces keep a size declaration local. Prefer the document's normal styles for
reports rather than manually resizing each paragraph.
Use \verb|\begin{enumerate}| with \verb|\item| for a numbered list.
For a two-column section use \verb|\twocolumn| and later \verb|\onecolumn|;
these commands normally start a new page, so use them deliberately.

\section{Equations and references}
\begin{lstlisting}
Inline mathematics: $F=ma$.
\begin{equation}
  T = 2\pi\sqrt{\frac{L}{g}}
  \label{eq:period}
\end{equation}
See Equation~\ref{eq:period}.
\end{lstlisting}
Numbered equations can be referenced without manually entering their numbers.
Use \texttt{align} from \texttt{amsmath} for aligned multi-line work:
\begin{lstlisting}
\begin{align}
  F &= ma,\\
  E &= \frac{1}{2}mv^2.
\end{align}
\end{lstlisting}
Keep units upright, for example \verb|\(L=0.50\,\mathrm{m}\)|.
For advanced units learn \texttt{siunitx}; check package compatibility
before adding it to older templates.
State model assumptions, not just equations.

\subsection{More physics notation}
These examples preserve the original tutorial's wider mathematical coverage
using standard \texttt{amsmath}/\texttt{amssymb} notation; the
\texttt{physics} package is not required.
\begin{lstlisting}
E = mc^2
\left\langle \psi_i \mid \psi_j \right\rangle = \delta_{ij}
\mathbf{u}\cdot\mathbf{v} = \sum_{i=1}^{n}u_i v_i
\frac{dy}{dx}, \quad \frac{\partial f}{\partial x}
\int_0^1 x^2\,dx
\iiint_V f(x,y,z)\,dx\,dy\,dz
\begin{pmatrix} 1 & 2 \\ 3 & 4 \end{pmatrix}
\mathbf{F} = m\mathbf{a}
\nabla\cdot\mathbf{E} = \frac{\rho}{\epsilon_0}
c = 2.99792458\times10^8\,\mathrm{m\,s^{-1}}
\end{lstlisting}
Put these expressions inside maths delimiters. For example:
\[
 \left\langle \psi_i \mid \psi_j \right\rangle=\delta_{ij},
 \qquad
 \mathbf{u}\cdot\mathbf{v}=\sum_{i=1}^{n}u_i v_i,
 \qquad
 A=\begin{pmatrix}1&2\\3&4\end{pmatrix}.
\]
\[
 \frac{dy}{dx},\qquad
 \frac{\partial f}{\partial x},\qquad
 \int_0^1 x^2\,dx,\qquad
 \iiint_V f(x,y,z)\,dx\,dy\,dz.
\]
Use explicit bounds or a named integration domain. The small space
\verb|\,| before a differential improves readability.
For structured unit formatting, the optional
\href{https://ctan.org/pkg/siunitx}{siunitx} package provides \verb|\SI|.

\subsection{Source-code listings}
Use \texttt{listings} to typeset code rather than inserting a screenshot.
The following example demonstrates the source of a listing:
\begin{lstlisting}[escapechar=|]
\begin{lstlisting}[language=Python]
import numpy as np
x = np.linspace(0, 10, 100)
print(x)
|\textbackslash|end{lstlisting}
\end{lstlisting}
Follow your module's guidance about including analysis code in an appendix
or separate repository. Code shown in a listing is not executed by LaTeX.

\section{Figures without fighting floats}
First generate \texttt{fit\_result.pdf} using the
\href{../Example_Code/Exersize.ipynb}{Python notebook}. With
\verb|\usepackage{graphicx}| in the preamble:
\begin{lstlisting}
\begin{figure}[htbp]
  \centering
  \includegraphics{../Example_Code/fit_result.pdf}
  \caption{Synthetic position data and a linear model.
           Error bars show assumed position uncertainty.}
  \label{fig:fit}
\end{figure}
Figure~\ref{fig:fit} compares observations and model.
\end{lstlisting}
For Overleaf, upload the figure and change the path to
\texttt{fit\_result.pdf}. Prefer vector PDF for plots; use suitable raster
formats for photographs. The path is relative to this guide's directory;
filenames and case must match. In Overleaf, upload the figure and use its
project-relative path.

A figure is a \textbf{float}. \texttt{[htbp]} suggests placements, not an exact
position. Reference the figure in text; do not force every figure with
\texttt{[H]}. Put the label \emph{after} the caption.
Captions should explain quantities, units, error bars, and physical meaning.
See \href{https://www.overleaf.com/learn/latex/Inserting_Images}{Overleaf's
guide to figures and images}.

\subsection{An image example you can rebuild}
To practise including an external image, copy the optional
\texttt{LaTeX\_meme.png} beside your source (or upload it to Overleaf).
Then try this example:
\begin{lstlisting}
\begin{figure}[htbp]
  \centering
  \includegraphics[width=0.45\linewidth]{LaTeX_meme.png}
  \caption{An existing tutorial image used to demonstrate image inclusion.}
  \label{fig:tutorial-image}
\end{figure}
\end{lstlisting}
\begin{figure}[htbp]
  \centering
  \fbox{\parbox[c][30mm][c]{0.45\linewidth}{\centering Example figure}}
  \caption{A self-contained placeholder demonstrating figure sizing and references.}
  \label{fig:tutorial-image}
\end{figure}
Figure~\ref{fig:tutorial-image} uses a relative width so it scales with the
text area. The \texttt{[H]} placement option in the old tutorial requires
\texttt{float}; it prevents floating, so use it only when genuinely necessary.

\section{Tables and citations}
Use tables for numerical results, with units in headings. \texttt{booktabs}
provides \verb|\toprule|, \verb|\midrule|, and \verb|\bottomrule|. Round
results to sensible significant figures.

\begin{lstlisting}
\begin{table}[htbp]
  \centering
  \caption{Illustrative values for practising table layout.}
  \label{tab:example}
  \begin{tabular}{rr}
    \toprule
    $x$ & $y$ \\
    \midrule
    1 & 2 \\
    2 & 4 \\
    \bottomrule
  \end{tabular}
\end{table}
\end{lstlisting}

For a separate BibTeX bibliography, create \texttt{ref.bib} beside your
source with this entry:
\begin{lstlisting}
@misc{overleaf_examples,
  author = {{Overleaf}},
  title = {Overleaf LaTeX Examples},
  year = {2024},
  howpublished = {\url{https://www.overleaf.com/latex/examples}}
}
\end{lstlisting}
In your document, cite it and print the bibliography with:
\begin{lstlisting}
See the examples~\cite{overleaf_examples}.
\bibliographystyle{plain}
\bibliography{ref}
\end{lstlisting}
The citation key must match the BibTeX entry exactly. See
\href{https://www.overleaf.com/learn/latex/Bibliography_management_with_bibtex}{Overleaf's
BibTeX guide} for citation and bibliography syntax.
Check source metadata and whether the source supports your claim.
A formatted citation is not proof of correctness.
For a complete rendered citation example in this guide, see the
Overleaf examples collection~\cite{overleaf_examples}.
This guide instead uses an inline \texttt{thebibliography} environment, so
its source compiles without a separate \texttt{.bib} file or BibTeX.
Do not combine the two bibliography approaches in the same document.
For a clickable external link, use
\verb|\href{URL}{link text}|.

\section{A useful lab-report structure}
\begin{enumerate}
\item Aim and physical model, including assumptions.
\item Method, instrument calibration, and data collection.
\item Results with units and clear figures/tables.
\item Uncertainties, residuals, systematic effects, and limitations.
\item Conclusion supported by measurements.
\item References and appropriate supplementary analysis/code.
\end{enumerate}
Follow your module's rubric when it differs.

\section{Exercises and checks}
\begin{enumerate}
\item Change the title/author and rebuild; confirm the PDF changed.
\item Add a labelled equation and refer to it without typing its number.
\item Add a table row, clearly marking real versus synthetic data.
\item Insert the Python plot and write a meaningful caption.
\item Add a real source, cite it, and rebuild without question marks.
\item Reorder two figures and check that references still follow them.
\end{enumerate}

\section{Troubleshooting}
\begin{description}
\item[No PDF:] Read the first error, not just the last summary.
\item[Undefined control sequence:] Check spelling and required packages.
\item[Missing dollar inserted:] Check maths delimiters and ordinary-text underscores.
\item[Figure missing:] Check working directory, filename/case, and whether
the figure was generated.
\item[Citation question mark:] Check keys, bibliography name, BibTeX, and
two further LaTeX passes when using BibTeX. For this guide's inline
bibliography, rerun LaTeX to resolve citations.
\item[Figure moved:] This is normal float behaviour; use cross-references.
\item[Overfull box:] Inspect the PDF; shorten text/URLs or restructure maths.
\item[Missing package:] Use the TeX distribution manager or Overleaf;
an editor extension is not TeX itself.
\end{description}
Further reference:
\href{https://www.overleaf.com/learn}{Overleaf learning resources} and
\href{https://www.latex-project.org/help/documentation/}{LaTeX documentation}.
\begin{thebibliography}{1}
\bibitem{overleaf_examples}
Overleaf.
\emph{Overleaf LaTeX Examples}.
2024.
\url{https://www.overleaf.com/latex/examples}.
\end{thebibliography}
\end{document}
